\documentclass[a4paper, 10pt]{article}
\usepackage[a4paper,left=8mm,right=8mm,top=10mm,bottom=20mm]{geometry}
\usepackage[utf8]{inputenc}
\usepackage{kotex}
\usepackage{amsmath,amsfonts,amssymb,mathrsfs,mathtools}
\usepackage{adjustbox,float,graphicx,relsize}
\usepackage{xcolor}
\usepackage{tikz}
\usepackage{setspace}
\usepackage{expl3}
\usepackage{xparse}
\usepackage{enumitem}
\usepackage{forloop}
\usepackage{amsthm}
\usepackage{verbatim}
\usepackage{wrapfig}
\usepackage{lipsum}
\usepackage{fancyvrb}
\usepackage{balance}
\usepackage{mathpartir}
\usepackage{stmaryrd}
\usepackage{cprotect}
\usepackage{pbox}
\usepackage[normalem]{ulem}
\usepackage{dashbox}
\usepackage{soul}
\usepackage{xcolor}
\usepackage{multirow}
\usepackage{framed}
\usepackage{enumitem}
\usepackage{cleveref}
\crefname{lemma}{렘마}{렘마}

\title{프로그램분석 26F HW2}
\author{2024-21504 남연우(Yeonwoo Nam)}
\date{}

\makeatletter
\let\thetitle=\@title
\makeatother

\newtheorem{theorem}{정리}
\newtheorem{lemma}{렘마}

\begin{document}
% \newcommand{\ie}{\textit{i.e.,} }
% \newcommand{\cf}{\textit{cf.} }
% \newcommand{\eg}{\textit{e.g.,} }
% \newcommand{\etal}{\textit{et al.}}

\newcommand{\defeq}{\ensuremath{\triangleq}}

\newcommand{\ra}[0]{\rightarrow{}}
\newcommand{\Ra}[0]{\Rightarrow{}}
\newcommand{\la}[0]{\leftarrow{}}
\newcommand{\La}[0]{\Leftarrow{}}
\newcommand{\lra}[0]{\leftrightarrow{}}
\newcommand{\rla}[0]{\leftrightarrow{}}
\newcommand{\Lra}[0]{\Leftrightarrow{}}
\newcommand{\Rla}[0]{\Leftrightarrow{}}

\newcommand{\sem}[1]{{\llbracket #1 \rrbracket}}
% Concrete and abstract semantics, applied to a state or abstract state.
\newcommand{\csem}[2]{\sem{#1}(#2)}
\newcommand{\asem}[2]{\sem{#1}^{\sharp}(#2)}
\newcommand{\lfp}[1]{{l\!f\!p\ }}

\newcommand{\powset}[1]{\mathbb{P}(#1)}

\newcommand{\Bool}{\mathbb{B}}
\newcommand{\ite}[3]{#1 \;{?}\; #2 : #3}
\newcommand{\cbite}[3]{#1 \;{\bm{?}}\; #2 \bm{:} #3}
\newcommand{\btrue}{\texttt{true}}
\newcommand{\bfalse}{\texttt{false}}

\newcommand{\False}{\textrm{False}}

\newcommand{\Empty}{{\emptyset}}

\newcommand{\Unit}{\textrm{Unit}}
\newcommand{\unit}{\texttt{tt}}

\newcommand{\unl}[1]{{\underline{#1}}}
\newcommand{\semu}[1]{{\sem{\unl{#1}}}}

%%% Local Variables:
%%% mode: latex
%%% TeX-master: "paper"
%%% End:


\maketitle

\paragraph{Collecting Semantics.}
프로그램 $P$에 나타나는 변수들의 유한집합을 $V$라 하자.
변수의 주소와 정수값은 구별하여 다음과 같이 정의한다.
\[
 \mathrm{Val}=\mathbb Z\uplus\{\&x\mid x\in V\},\qquad
 \Sigma=V\to\mathrm{Val},\qquad D=\powset{\Sigma}
\]

좌변의 위치와 우변의 값은 각각 $\sem{L}_{\mathrm L}$과 $\sem{R}_{\mathrm R}$로 구별하여 정의한다.
$\sigma(x)\in\mathbb{Z}$일 때 $\sem{*x}$는 정의되지 않는다.
\[
\begin{aligned}
 \sem{x}_{\mathrm L}(\sigma)&=x,
 &\sem{n}_{\mathrm R}(\sigma)&=n,\\
 \sem{*x}_{\mathrm L}(\sigma)&=y\quad(\sigma(x)=\&y),
 &\sem{x}_{\mathrm R}(\sigma)&=\sigma(x),\\
 &&\sem{*x}_{\mathrm R}(\sigma)&=\sigma(y)\quad(\sigma(x)=\&y),\\
 &&\sem{\&x}_{\mathrm R}(\sigma)&=\&x.
\end{aligned}
\]
우변과 좌변의 위치는 모두 대입 전 상태에서 계산한다.
$B$는 메모리를 변경하지 않는 Boolean 식으로 둔다.

실행 중 나타나는 상태를 over-approximate하는 collecting semantics $\sem{C}\in D\ra D$를 정의한다.
종료 상태와 중간 상태를 구별하지 않고, sequence에서는 앞 statement가 모은 모든 상태를
뒷 statement의 입력으로 사용한다. 따라서 실제로 실행되지 않는 순서로 얻은 상태도 포함할 수 있다.
분석의 목적은 모든 실제 points-to 관계를 포함하는 것이므로 이 over-approximate를 허용한다.
Boolean 식에 의한 state 집합의 제한을
\[
 \mathcal B_{\btrue}(S)=\{\sigma\in S\mid\sem{B}(\sigma)=\btrue\}
\]
로 두고, statement의 의미를 다음과 같이 정의한다.
\[
\begin{aligned}
 \csem{L:=R}{S}
  &=S\cup\{\sigma[u\mapsto v]\mid\sigma\in S,\
        \sem{L}_{\mathrm L}(\sigma)=u,\
        \sem{R}_{\mathrm R}(\sigma)=v\},\\
 \csem{C_1;C_2}{S}&=\csem{C_2}{\csem{C_1}{S}},\\
 \csem{\mathsf{while}\ B\ C}{S}
  &=\operatorname{lfp}\bigl(\lambda X.~S\cup\csem{C}{\mathcal B_{\btrue}(X)}\bigr).
\end{aligned}
\]
대입에서 좌변이나 우변의 계산이 정의되지 않으면 그 상태에서 새 상태를 만들지 않는다.
각 함수는 포함관계에 대해 단조이며 합집합을 보존하고, $S\subseteq\csem{C}{S}$이다.
루프의 의미는 $X_0=\emptyset$,
$X_{k+1}=S\cup\csem{C}{\mathcal B_{\btrue}(X_k)}$에 대한
$\bigcup_{k\ge0}X_k$이다.
본문의 중간 상태도 다음 반복의 입력 후보로 사용하므로 이 역시 과over-approximate이다.

이 정의는 모든 실제 실행 중 나타나는 상태를 포함한다.
초기에 points-to 관계가 없다고 가정하고, 초기상태 집합을
\[
 I_0=\{\sigma\in\Sigma\mid\forall x\in V.~\sigma(x)\in\mathbb Z\}
\]
로 고정한다. 따라서 모든 변수의 초기값은 임의의 정수이다.
이하에서는 $\csem{P}{I_0}$를 요약의미가 포함함을 보인다.

\paragraph{Abstract Domain.}
요약공간을 다음과 같이 정의한다. $A(x)$는 $x$가 가리킬 \emph{수 있는} 변수들의 집합이다.
\[
 A\in D^\sharp=V\ra\powset{V}.
\]
\[
\begin{gathered}
 A\sqsubseteq A'\ \Longleftrightarrow\
       \forall x\in V.\ A(x)\subseteq A'(x),\qquad
 \bot^\sharp(x)=\emptyset,\quad\top^\sharp(x)=V,\\
 \left(\bigsqcup_i A_i\right)(x)=\bigcup_i A_i(x),
 \qquad
 \left(\bigsqcap_i A_i\right)(x)=\bigcap_i A_i(x).
\end{gathered}
\]
따라서 $D^\sharp$는 finite complete lattice이다. abstract/concretization 함수를
\[
 \alpha(S)(x)=\{y\in V\mid\exists\sigma\in S.\ \sigma(x)=\&y\},
 \qquad
 \gamma(A)=\{\sigma\in\Sigma\mid
    \forall x,y\in V.\ \sigma(x)=\&y\Rightarrow y\in A(x)\}
\]
로 정의한다. 그러면
\[
 \alpha(S)\sqsubseteq A\quad\Longleftrightarrow\quad S\subseteq\gamma(A)
\]
이므로 $(\alpha,\gamma)$는 Galois connection을 이룬다.
정의에 의해 $I_0=\gamma(\bot^\sharp)$이고 $\alpha(I_0)=\bot^\sharp$이다.
따라서 실제 도달 상태 집합을 $\mathrm{Reach}_P(I_0)$라 쓰면, 증명할 포함관계는
\[
 \mathrm{Reach}_P(I_0)\subseteq\csem{P}{I_0}
 \subseteq\gamma(\asem{P}{\bot^\sharp})
\]
이다.

\newpage
\paragraph{Abstract Semantics.}
좌변이 쓸 수 있는 위치와 우변이 가질 수 있는 \emph{주소}를 각각
\[
\begin{aligned}
 \sem{x}^{\sharp}_{\mathrm L}(A)&=\{x\},
 &\sem{n}^{\sharp}_{\mathrm R}(A)&=\emptyset,\\
 \sem{*x}^{\sharp}_{\mathrm L}(A)&=A(x),
 &\sem{x}^{\sharp}_{\mathrm R}(A)&=A(x),\\
 &&\sem{*x}^{\sharp}_{\mathrm R}(A)&=\bigcup_{y\in A(x)}A(y),\\
 &&\sem{\&x}^{\sharp}_{\mathrm R}(A)&=\{x\}.
\end{aligned}
\]
로 정의한다. Abstract semantics $\sem{C}^\sharp:D^\sharp\to D^\sharp$를 아래와 같이 정의한다.
\[
\begin{aligned}
 \asem{L:=R}{A}(u)&=A(u)\cup
 \begin{cases}
  \sem{R}^{\sharp}_{\mathrm R}(A)&u\in\sem{L}^{\sharp}_{\mathrm L}(A),\\
  \emptyset&\text{그 외},
 \end{cases}\\
 \asem{C_1;C_2}{A}&=\asem{C_2}{\asem{C_1}{A}},\\
 \asem{\mathsf{while}\ B\ C}{A}
  &=\operatorname{lfp}\bigl(\lambda X.~A\sqcup\asem{C}{X}\bigr).
\end{aligned}
\]
요약공간은 정수값이나 조건식의 참/거짓을 기록하지 않으므로,
조건에 대한 요약함수는 $\mathcal B^\sharp_{\btrue}(A)=A$로 둔다.
실제로 $S\subseteq\gamma(A)$이면
$\mathcal B_{\btrue}(S)\subseteq S\subseteq\gamma(A)$이므로 이 구현은 안전하다.
루프의 정의에서는 이 항등함수를 생략하였다.

$A\sqsubseteq\asem{C}{A}$이다.
루프의 최소고정점은
\[
 X^{\sharp}_0=\bot^\sharp,\qquad
 X^{\sharp}_{k+1}=A\sqcup\asem{C}{X^{\sharp}_k}
\]
로 계산한다. $D^\sharp$가 유한하므로
유한 번의 반복 후 안정되며 widening은 필요 없다.
초기 요약상태는 $\alpha(I_0)=\bot^\sharp$이며,
프로그램 전체의 분석 결과는 $\asem{P}{\bot^\sharp}$이다.


\paragraph{Soundness Proof.}
\begin{lemma}
$S\subseteq\gamma(A)$이면 $\csem{L:=R}{S}\subseteq\gamma(\asem{L:=R}{A})$이다.
\end{lemma}
\begin{proof}
$\sigma\in S$를 잡으면 $\sigma\in\gamma(A)$이다.

먼저 $\sem{L}_{\mathrm L}(\sigma)\in\sem{L}^{\sharp}_{\mathrm L}(A)$임을 증명하자.
$L=x$이면 자명하고, $L=*x$이면 $\sigma(x)=\&y$와 $\sigma\in\gamma(A)$로부터 $y\in A(x)$이다.

또 $\sem{R}_{\mathrm R}(\sigma)=\&v$이면 $v\in\sem{R}^{\sharp}_{\mathrm R}(A)$임을 증명할 수 있다.
$R=n$인 경우는 불가능하고, $R=\&x$이면 $v=x$이다.
$R=x$이면 $\gamma$의 정의에 의해 참이다.
$R=*x$이면 어떤 $y$에 대해 $\sigma(x)=\&y$, $\sigma(y)=\&v$이므로
$y\in A(x)$, $v\in A(y)$이다.

$s=(L:=R)$라 하고 $\sigma'\in\csem{s}{S}$를 잡자.
$\sigma'\in S$이면 $A\sqsubseteq \asem{s}{A}$로부터 원하는 결론을 얻는다.
그 외에는 어떤 $\sigma\in S$와 $u,v$에 대해
$\sem{L}_{\mathrm L}(\sigma)=u$, $\sem{R}_{\mathrm R}(\sigma)=v$이고
$\sigma'=\sigma[u\mapsto v]$이다.
$\sigma'(a)=\&b$인 임의의 $a,b$를 잡는다.
$a\ne u$이면 $\sigma(a)=\&b$이므로 $b\in A(a)\subseteq \asem{s}{A}(a)$이다.
$a=u$이면 우변의 값이 $\&b$이므로 위 성질에 따라
$b\in\sem{R}^{\sharp}_{\mathrm R}(A)\subseteq \asem{s}{A}(u)$이다.
따라서 $\sigma'\in\gamma(\asem{s}{A})$이다.
\end{proof}

\begin{theorem}
$S\in D$, $A\in D^\sharp$, statement $C$에 대해
\[
 S\subseteq\gamma(A)\quad\Longrightarrow\quad
 \csem{C}{S}\subseteq\gamma(\asem{C}{A}).
\]
특히 $S=I_0$, $A=\bot^\sharp$로 두면
$\alpha(\csem{P}{I_0})\sqsubseteq\asem{P}{\bot^\sharp}$이다.
\end{theorem}
\begin{proof}
$C$의 구조에 대해 induction으로 증명한다. $C=(L:=R)$인 경우는 렘마 1에 의해 성립한다.

$C = (C_1;C_2)$이면 induction hypothesis를 차례로 적용하여
\[
 \csem{C_1}{S}\subseteq\gamma(\asem{C_1}{A}),\qquad
 \csem{C_2}{\csem{C_1}{S}}
 \subseteq\gamma(\asem{C_2}{\asem{C_1}{A}})
\]
를 얻으므로 정의에 의해 결론이 성립한다.

$C=W=\mathsf{while}\ B\ C_0$인 경우 다음 두 함수를 정의한다.
\[
 F(X)=S\cup\csem{C_0}{\mathcal B_{\btrue}(X)},\qquad
 F^\sharp(Y)=A\sqcup\asem{C_0}{Y}.
\]
$F$는 연속이고 $F^\sharp$는 단조이다.
임의의 $Y\in D^\sharp$에 대해
$\mathcal B_{\btrue}(\gamma(Y))\subseteq\gamma(Y)$이므로,
$C_0$에 대한 induction hypothesis와 $S\subseteq\gamma(A)$로부터
\[
\begin{aligned}
 F(\gamma(Y))
 &=S\cup\csem{C_0}{\mathcal B_{\btrue}(\gamma(Y))}\\
 &\subseteq\gamma(A)\cup\gamma(\asem{C_0}{Y})\\
 &\subseteq\gamma\bigl(A\sqcup\asem{C_0}{Y}\bigr)
 =\gamma(F^\sharp(Y)).
\end{aligned}
\]
따라서 $F\circ\gamma\sqsubseteq\gamma\circ F^\sharp$이다.
Fixpoint Transfer Theorem을 적용하면
\[
 \csem{W}{S}=\operatorname{lfp}(F)
 \subseteq\gamma\left(\bigsqcup_{i\ge0}(F^\sharp)^i(\bot^\sharp)\right)
 =\gamma(\operatorname{lfp}(F^\sharp))
 =\gamma(\asem{W}{A}).
\]

마지막 주장은 $I_0=\gamma(\bot^\sharp)$와 Galois connection으로부터 따른다.
\end{proof}
이 정리와 $\mathrm{Reach}_P(I_0)\subseteq\csem{P}{I_0}$를 결합하면,
모든 실제 실행 중 나타나는 points-to 관계가 분석 결과에 포함됨을 증명한다.

\end{document}